% Research Paper Template (Modern Best Practices)
% Compile: latexmk -pdf main.tex
% Or: pdflatex → biber → pdflatex → pdflatex
\documentclass[11pt,a4paper]{article}

%% ===========================================
%% ENCODING AND FONTS
%% ===========================================
\usepackage[utf8]{inputenc}
\usepackage[T1]{fontenc}
\usepackage{lmodern}           % Modern, fully embeddable fonts
\usepackage{microtype}         % Improved typography

%% ===========================================
%% PAGE LAYOUT
%% ===========================================
\usepackage[margin=1in]{geometry}
\usepackage{setspace}
\onehalfspacing                % 1.5 line spacing (adjust as needed)

%% ===========================================
%% MATHEMATICS
%% ===========================================
\usepackage{amsmath,amssymb,amsthm}
\usepackage{mathtools}

% Theorem environments
\newtheorem{theorem}{Theorem}[section]
\newtheorem{lemma}[theorem]{Lemma}
\newtheorem{proposition}[theorem]{Proposition}
\newtheorem{corollary}[theorem]{Corollary}
\theoremstyle{definition}
\newtheorem{definition}{Definition}[section]
\newtheorem{example}{Example}[section]
\theoremstyle{remark}
\newtheorem{remark}{Remark}[section]

%% ===========================================
%% GRAPHICS AND FIGURES
%% ===========================================
\usepackage{graphicx}
\graphicspath{{figures/}}      % Default figures directory
\usepackage{subcaption}        % Subfigures

%% ===========================================
%% TABLES
%% ===========================================
\usepackage{booktabs}          % Professional rules
\usepackage{siunitx}           % Number alignment and units
\sisetup{
    round-mode = places,
    round-precision = 2,
    detect-weight = true
}
\usepackage{tabularx}          % Auto-width columns

%% ===========================================
%% CODE LISTINGS (choose one)
%% ===========================================
% Option 1: minted (requires -shell-escape)
% \usepackage{minted}

% Option 2: listings (no shell-escape needed)
\usepackage{listings}
\usepackage{xcolor}
\lstset{
    basicstyle=\ttfamily\small,
    keywordstyle=\color{blue},
    commentstyle=\color{gray},
    numbers=left,
    numberstyle=\tiny\color{gray},
    breaklines=true,
    frame=single
}

%% ===========================================
%% BIBLIOGRAPHY (Modern - biblatex + biber)
%% ===========================================
\usepackage[
    backend=biber,
    style=numeric,             % Or: authoryear, alphabetic, apa, ieee
    sorting=nyt,               % Name, year, title
    maxbibnames=99
]{biblatex}
\addbibresource{references.bib}

%% ===========================================
%% CROSS-REFERENCES (load order critical!)
%% ===========================================
\usepackage{hyperref}
\hypersetup{
    colorlinks=true,
    linkcolor=blue,
    citecolor=blue,
    urlcolor=cyan,
    pdftitle={Your Paper Title},
    pdfauthor={Your Name},
    pdfkeywords={keyword1, keyword2, keyword3}
}
\usepackage[noabbrev,capitalise]{cleveref}  % MUST be after hyperref

%% ===========================================
%% DOCUMENT METADATA
%% ===========================================
\title{Your Paper Title}
\author{
    First Author\thanks{Affiliation, \texttt{email@example.com}}
    \and
    Second Author\thanks{Affiliation, \texttt{email@example.com}}
}
\date{\today}

%% ===========================================
%% DOCUMENT
%% ===========================================
\begin{document}

\maketitle

\begin{abstract}
Write your abstract here. It should be a concise summary of your research,
typically 150--300 words, covering the problem, methods, key results, and conclusions.
Use an en-dash (--) for ranges, not a hyphen.
\end{abstract}

\noindent\textbf{Keywords:} keyword1, keyword2, keyword3

\section{Introduction}
\label{sec:intro}

Introduce your research problem and motivation. When citing, use
\textcite{example_key} for ``Author (Year)'' or \cite{example_key} for ``[1]''.

Use smart cross-references: see \cref{sec:methods} for methodology and
\cref{fig:example} for results visualization.

\section{Related Work}
\label{sec:related}

Discuss previous research. Multiple citations: \cite{example_key,book_example}.

\section{Methodology}
\label{sec:methods}

Describe your approach. Include equations with labels:
\begin{equation}
    \mathcal{L}(\theta) = -\sum_{i=1}^{N} y_i \log \hat{y}_i + (1-y_i) \log(1-\hat{y}_i)
    \label{eq:loss}
\end{equation}

Reference equations with \cref{eq:loss}. For multiple: \cref{eq:loss,eq:precision}.

\begin{align}
    \text{Precision} &= \frac{\text{TP}}{\text{TP} + \text{FP}} \label{eq:precision} \\
    \text{Recall} &= \frac{\text{TP}}{\text{TP} + \text{FN}} \label{eq:recall}
\end{align}

\section{Results}
\label{sec:results}

Present your findings with figures and tables.

% Example figure
\begin{figure}[htbp]
    \centering
    % \includegraphics[width=0.8\textwidth]{your_figure.pdf}
    \fbox{\parbox{0.8\textwidth}{\centering\vspace{2cm}Figure Placeholder\vspace{2cm}}}
    \caption{Description of the figure. Explain what the reader should observe.}
    \label{fig:example}
\end{figure}

% Example table with siunitx alignment
\begin{table}[htbp]
    \centering
    \caption{Model performance comparison. Best results in bold.}
    \label{tbl:results}
    \sisetup{table-format=2.2}
    \begin{tabular}{l S S[table-format=1.3] S[table-format=2.1]}
        \toprule
        Model & {Accuracy (\%)} & {F1 Score} & {Latency (ms)} \\
        \midrule
        Baseline    & 75.32 & 0.723 & 12.4 \\
        ResNet-50   & 89.15 & 0.871 & 45.2 \\
        Transformer & 92.48 & 0.912 & 78.3 \\
        \textbf{Ours} & \bfseries 94.21 & \bfseries 0.934 & 52.1 \\
        \bottomrule
    \end{tabular}
\end{table}

As shown in \cref{tbl:results}, our method achieves the best performance.

\section{Discussion}
\label{sec:discussion}

Interpret results and discuss implications. Address limitations.

\section{Conclusion}
\label{sec:conclusion}

Summarize contributions:
\begin{enumerate}
    \item First contribution
    \item Second contribution
    \item Third contribution
\end{enumerate}

Future work includes\ldots

\section*{Acknowledgments}

This work was supported by [funding source]. We thank [collaborators].

%% ===========================================
%% BIBLIOGRAPHY
%% ===========================================
\printbibliography

%% ===========================================
%% APPENDIX (optional)
%% ===========================================
\appendix
\section{Additional Results}
\label{app:additional}

Supplementary material referenced from main text.

\end{document}
